\section{Problem Definition and Use Cases}

\subsection{Problem Definition}

Modern children face severe attention span challenges that make traditional programming education ineffective. 
Standard programming courses require prolonged focus, abstract thinking, and patience for delayed results. 
These requirements clash with how children consume information today. 
Minecraft offers a solution because children already invest hours in the game and want to improve their gameplay. 
Current Minecraft programming tools demand too much technical knowledge and fail to provide immediate, visible outcomes. 
Young learners need a programming environment where code produces instant visual results within a game they love. 
This gap creates an opportunity for a domain-specific language that turns programming into an engaging experience.

\subsection{Technical Barriers in Existing Tools}

Children face Java requirements far above beginner level when working with Minecraft mods. 
Tools like ComputerCraft use Lua syntax that confuses young learners. 
Setup complexity takes time before anyone writes a single line of code, which stops learning early. 
Command blocks, an in-game programming interface, overwhelm beginners from the start. 
Parents or teachers must handle installation and configuration because children lack the technical knowledge.

\subsection{Pedagogical Gaps}

Children learn abstract concepts without seeing corresponding actions in the game. 
Block-based tools hide text syntax and create barriers when transitioning to real code. 
Motivation drops when results feel distant or unclear. 
Available tools do not match how children interact with digital products today. 
Educational resources assume adult knowledge or prior programming experience.

\subsection{User Experience Problems}

Text output appears with no connection to in-game action. 
Feedback loops take too long after writing code. 
Error messages confuse rather than guide young learners. 
No clear difficulty progression exists to match skill growth. 
Interface complexity becomes an obstacle instead of supporting the learning process.

\subsection{Specific Problems the DSL Solves}

\begin{itemize}
    \item[\textbf{--}] \textbf{Lack of Immediate Visual Feedback:} 
    Children write code and expect to see action happen in the game world. 
    Existing tools separate code execution from visible results. 
    Learning slows when outcomes stay invisible or appear only as text in a console. 
    The DSL runs commands directly inside the game environment. 
    Children watch their bot perform actions the moment they execute the script. 
    Blocks appear, items move, and structures are built in real time while the code runs.

    \item[\textbf{--}] \textbf{Weak Path Toward Real Programming:} 
    Children learn logic through drag-and-drop blocks, then face text-based code later without preparation. 
    The transition feels abrupt and confusing. 
    The DSL uses simple text commands from the beginning. 
    Children read and write actual code syntax instead of clicking visual blocks. 
    The syntax remains clear and readable while teaching real programming structure. 
    When children move to other languages later, they already understand how text-based code works.

    \item[\textbf{--}] \textbf{Excessive Technical Complexity:} 
    Current tools require deep technical knowledge to accomplish basic tasks. 
    Setting up mods, understanding game mechanics, and navigating complex interfaces create a steep learning curve. 
    The DSL hides these complex mechanics behind clear, simple commands. 
    Children focus on programming logic and problem structure instead of fighting technical barriers. 
    Commands abstract away implementation details while still teaching fundamental concepts.

    \item[\textbf{--}] \textbf{Poor Alignment with Modern Learning Patterns:} 
    Traditional programming education uses long tutorials with delayed gratification. 
    Children lose interest before seeing meaningful results. 
    The DSL embeds learning directly inside Minecraft, where children already spend time. 
    Each command triggers a visible change in the game world. 
    Engagement stays high because every action produces immediate outcomes that children care about.

    \item[\textbf{--}] \textbf{Missing Support for Independent Learning:} 
    Children depend on adult help to progress through existing programming tools. 
    Exploration stops when guidance becomes unavailable. 
    The DSL supports independent trial and error through clear syntax and fast feedback loops. 
    Children experiment with commands, see what happens, fix mistakes when things go wrong, and advance on their own. 
    Parents without programming knowledge no longer need to intervene constantly.

    \item[\textbf{--}] \textbf{Lack of Meaningful Context:} 
    Programming exercises in traditional curricula feel arbitrary and disconnected from student interests. 
    Children solve problems they do not care about, so motivation fades quickly. 
    The DSL ties every coding task to goals children set for themselves in Minecraft. 
    Building a house, automating a farm, or organizing inventory all serve purposes that children choose. 
    Each line of code accomplishes something personally meaningful.
\end{itemize}

\subsection{Use Case Scenarios}

\begin{itemize}
    \item[\textbf{--}] \textbf{Build a House:}
    The goal is for the user to place blocks in a consistent pattern to form a simple structure. 
    A sequence of movement and placement commands is written to define where each block goes. 
    Correct results depend on executing commands in the right order, and repeating similar actions helps reinforce the idea of structured sequences and repetition.

    \item[\textbf{--}] \textbf{Create a Farming Bot:}
    The objective is to plant crops in a straight line efficiently. 
    A repeat instruction combines movement and planting actions so the same steps are applied across multiple blocks. 
    This use case introduces loops and demonstrates how automation reduces manual effort.

    \item[\textbf{--}] \textbf{Manage Inventory:}
    The goal is to remove or move items from specific inventory slots. 
    Inventory and drop commands are written with parameters such as slot numbers and coordinates. 
    This use case teaches how structured commands and parameters control precise behavior.

    \item[\textbf{--}] \textbf{React to Conditions:}
    The objective is for the character to respond to changes in game state, such as eating when hunger drops below a certain level. 
    A condition is defined with an associated action. 
    This use case introduces decision-making and conditional logic.

    \item[\textbf{--}] \textbf{Chest Management Bot:}
    The goal is to organize collected items into storage chests by type. 
    A bot is programmed to move through a base, identify items, and place them into the correct chests. 
    This use case combines movement, logic, and automation, reinforcing data organization and control flow through a practical task.
\end{itemize}

\subsection{Pain Points in Existing Solutions for Teaching Programming}

\begin{itemize}
    \item[\textbf{--}] \textbf{Scratch and Other Visual Languages:}
    Scratch lowers the entry barrier by removing syntax errors, but it operates in an abstract environment disconnected from the games children already play. 
    The visual block system hides real text-based programming syntax, which makes the transition to textual languages like Python or Java difficult \cite{noauthor_connecting_2024}. 
    As projects grow in complexity, block-based coding limits the types of programs learners can build \cite{guide_scratch_2023}. 
    Learners often struggle to connect their work to familiar gaming environments, causing engagement to drop \cite{adler_improving_nodate}.

    \item[\textbf{--}] \textbf{Minecraft Education Edition:}
    Minecraft Education Edition integrates learning into a familiar game, but access requires institutional licensing, limiting home use \cite{noauthor_licensing_nodate}. 
    The platform targets classroom instruction and emphasizes guided lessons over open-ended experimentation \cite{noauthor_pros_2025}. 
    Activities rely heavily on pre-built curricula, which reduces creative freedom and limits exploration. 
    Many children prefer the regular version of Minecraft where friends play, making the educational edition feel isolated from their social gaming environment \cite{sripan_gamified_2025}.

    \item[\textbf{--}] \textbf{ComputerCraft:}
    ComputerCraft introduces real programming concepts through Lua, but it requires text-based coding from the start. 
    Installation and setup demand technical knowledge often unavailable to parents. 
    Documentation assumes prior programming experience, which can overwhelm beginners. 
    Young learners face high complexity before achieving results, leading to frustration and early disengagement \cite{kazemitabaar_scaffolding_2023}.

    \item[\textbf{--}] \textbf{Other Common Approaches and Their Limitations:}
    Minecraft command blocks allow in-game automation, but dense syntax produces frequent errors that confuse beginners. 
    Redstone teaches logic indirectly, and debugging circuits is unclear, making reasoning about failures difficult. 
    Traditional programming languages such as Python or Java require setup and configuration before learners can see results; syntax errors in early stages often block progress and reduce motivation.
\end{itemize}

\subsection{Alternative DSLs}

\begin{itemize}
    \item[\textbf{--}] \textbf{DSL for Modding Minecraft:}
    DSL for Modding Minecraft is a language designed to help users write Minecraft mods at a high level by generating Java code behind the scenes. 
    It abstracts away the complexity of setting up and coding in Java, focusing on domain concepts like blocks, items, and entities. 
    This DSL targets mod creation rather than scripting character behavior inside the game, and requires code generation and build tools to produce an installable mod file \cite{stalla_dsl_2024}.

    \item[\textbf{--}] \textbf{Minescript:}
    Minescript is a mod that lets users write Python scripts to control Minecraft gameplay. 
    Scripts run from the in-game console and can access game state, place blocks, and interact with the world. 
    Although it uses a real language (Python) rather than a specialized DSL syntax, its integration with Minecraft and simplified access to game functions make it a scripting alternative for controlling actions programmatically within the game world \cite{noauthor_minescript_nodate}.

    \item[\textbf{--}] \textbf{SpigotDSL:}
    SpigotDSL is a Kotlin-based domain-specific library that makes writing server plugins more readable. 
    It leverages Kotlin's DSL features to express plugin logic (such as event handlers) in a concise way compared to raw Java. 
    It is a DSL in the sense of domain-focused language design, but it is tied to server plugin development rather than in-game behavior scripting for learners \cite{noauthor_spigotdsl_nodate}.
\end{itemize}

\subsection{Problem statement }

Children disengage from programming education when learning tools fail to provide immediate visual feedback and do not connect abstract coding concepts to activities that sustain their interest.
